%% kapa: exam=1 round=34 year=2023 period=1 form=A subject=economics subject_ko=경제학원론 seq=4 number=44
%% source: Q-Net 감정평가사 기출문제 (cst00309) artlSeq=5211789 file=raw/1008/5211789/2023년도 제34회 감정평가사 1차시험 1교시.pdf
%% conversion: kordoc 4.24.51 -> tools/gen1.py (auto); review=proofread-sonnet; flags=none
\begin{question}{44}
두 재화 $X$, $Y$를 소비하는 갑의 효용함수가 $U = XY^2$이고, $X$재의 가격은 1, $Y$재의 가격은 2, 소득은 90이다. 효용함수와 소득이 각각 $U = \sqrt{XY}$, 100으로 변경되었을 경우, 갑의 효용이 극대화되는 $X$재와 $Y$재의 구매량의 변화로 옳은 것은?
\begin{choices}
  \item $X$재 10 증가, $Y$재 5 감소
  \item $X$재 10 증가, $Y$재 5 증가
  \item $X$재 20 증가, $Y$재 5 감소
  \item $X$재 20 증가, $Y$재 10 감소
  \item $X$재 20 증가, $Y$재 10 증가
\end{choices}
\end{question}
\answer{3}
